\documentclass[../main.tex]{subfiles}
\begin{document}

\section{Summary}

Sinh viên nhắc lại các vấn đề mà đồ án đã giải quyết được, cũng như những vấn đề còn tồn đọng của đồ án.  

\section{Suggestion for Future Works }

Sinh viên đề xuất hướng phát triển trong tương lai (nếu có) .

\subsection{Giới hạn và tính bảo mật của IDS}

IDS là thành phần quan trọng nhưng không phải giải pháp ``all-in-one'' trong hạ tầng mạng. Một NIDS dù mạnh vẫn bị giới hạn bởi ba yếu tố:

\begin{itemize}
    \item \textbf{Tầm nhìn}: IDS có thực sự quan sát hết lưu lượng không? Kẻ tấn công có thể định tuyến đi vòng qua IDS.
    \item \textbf{Tải trọng (hiệu suất)}: khi lưu lượng quá lớn (ví dụ tấn công DoS), IDS có thể bị ``ngạt'' và bỏ lỡ gói độc hại.
    \item \textbf{Độ sâu phân tích}: phân tích càng sâu càng tốn tài nguyên; parser kém hiệu quả trở thành nút thắt, trong khi phân tích nông có thể bỏ sót tấn công tinh vi.
\end{itemize}

\textbf{Lỗi triển khai}: IDS là phần mềm phức tạp và cũng có lỗ hổng. Parser lỗi có thể khiến IDS treo khi nhận gói tin đặc biệt. IDS thường viết bằng C/C++ để đạt tốc độ cao nên dễ gặp tràn bộ đệm (\textit{buffer overflow}); nếu bị chiếm quyền, kẻ tấn công có vị trí thuận lợi trong mạng. Các biện pháp giảm thiểu gồm:
\begin{itemize}
    \item Dùng ngôn ngữ an toàn bộ nhớ (\textit{memory-safe languages}). Thay vì C/C++, có thể chuyển sang Rust --- Suricata đang chuyển dần phần phân tích giao thức sang Rust để giảm lỗi tràn bộ nhớ.
    \item \textbf{Fuzzing}: đưa dữ liệu lỗi vào IDS để phát hiện lỗi trước kẻ tấn công khai thác.
    \item \textbf{Sanitizing}: phân tích hành vi để ngăn thực thi lệnh trái phép từ dữ liệu đầu vào.
\end{itemize}

\end{document}